\documentclass{mcom-l}

% Essential packages for the content
\usepackage{amssymb}
\usepackage{graphicx}
\usepackage{url}
\usepackage{algorithm}
\usepackage{algpseudocode}
\usepackage{longtable}
\usepackage{booktabs}
\usepackage{array}
\usepackage{pgfplots}

\pgfplotsset{compat=1.18}

\newtheorem{theorem}{Theorem}[section]
\newtheorem{lemma}[theorem]{Lemma}
\newtheorem{corollary}[theorem]{Corollary}
\newtheorem{proposition}[theorem]{Proposition}

\theoremstyle{definition}
\newtheorem{definition}[theorem]{Definition}
\newtheorem{example}[theorem]{Example}
\newtheorem{xca}[theorem]{Exercise}

\theoremstyle{remark}
\newtheorem{remark}[theorem]{Remark}
\newtheorem{assumption}[theorem]{Assumption}

\numberwithin{equation}{section}

%    Absolute value notation
\newcommand{\abs}[1]{\lvert#1\rvert}

%    Blank box placeholder for figures (to avoid requiring any
%    particular graphics capabilities for printing this document).
\newcommand{\blankbox}[2]{%
  \parbox{\columnwidth}{\centering
%    Set fboxsep to 0 so that the actual size of the box will match the
%    given measurements more closely.
    \setlength{\fboxsep}{0pt}%
    \fbox{\raisebox{0pt}[#2]{\hspace{#1}}}%
  }%
}

% Algorithm environment settings
\renewcommand{\algorithmicrequire}{\textbf{Input:}}
\renewcommand{\algorithmicensure}{\textbf{Output:}}

\begin{document}
\raggedbottom

\title[Four cubes up to \(10^9\)]{Verifying the Sum of Four Cubes Problem up to \(10^9\) via Generalized Pell Equations}

%    Information for first author
\author{Kaiyi Zhang}
\address{Institute for Advanced Study, Tsinghua University, Beijing, China}
\email{kaiyizhang@tsinghua.edu.cn}

\author{Zhuo Huang}
\address{School of Computer Science, Shanghai Jiao Tong University, Shanghai, China}
\email{sh1kaku@sjtu.edu.cn}

\author{Anyu Wang}
\address{Institute for Advanced Study, Tsinghua University, Beijing, China; State Key Laboratory of Cryptography and Digital Economy Security, Tsinghua University, Beijing, China; School of Cryptologic Science and Engineering, Shandong University, Jinan, China; Zhongguancun Laboratory, Beijing, China; National Financial Cryptography Research Center, Beijing, China}
\email{anyuwang@tsinghua.edu.cn}
\thanks{Corresponding author: Anyu Wang}

\author{Xiaoyun Wang}
\address{Institute for Advanced Study, Tsinghua University, Beijing, China; State Key Laboratory of Cryptography and Digital Economy Security, Tsinghua University, Beijing, China; School of Cryptologic Science and Engineering, Shandong University, Jinan, China; Zhongguancun Laboratory, Beijing, China; National Financial Cryptography Research Center, Beijing, China; Shandong Institute of Blockchain, Shandong, China}
\email{xiaoyunwang@tsinghua.edu.cn}
%    General info
\subjclass[2020]{Primary 11Y50; Secondary 11D25, 11P05, 11Y16}

\date{}
%\dedicatory{This paper is dedicated to our advisors.}

\keywords{Sum of four cubes, generalized Pell equations, computational verification, Diophantine equations, computational number theory}

\begin{abstract}
We computationally verify the sum of four cubes problem for every positive
integer $n \leq 10^9$, producing explicit integer representations
$n=x_1^3+x_2^3+x_3^3+x_4^3$.  The search is based on a transformation of a
structured four-cube identity into congruence-constrained generalized Pell
equations $u^2-d v^2=M$.  For each candidate representation,
correctness is checked by exact integer arithmetic, and the associated code
and data are made publicly available.  We emphasize the computation as a
certificate-based verification: the final checking pass is independent of
the Pell search and verifies only coverage and the four-cube identities.
We also give a local-admissibility heuristic, based on the congruence
obstructions forced by primes dividing $d$ and on conditional
Bernays-type density, which explains the observed scale of the search
parameter.  The same method also finds representations for large integers,
such as $n=9\cdot 2^{128}+5$, within seconds.
\end{abstract}

\maketitle

\section{Introduction}
\label{sec:introduction}

\subsection{Problem background}
\label{subsec:problem_background}

The sum of four cubes problem asks whether every integer $n$ can be
expressed as the sum of four integer cubes, i.e., whether there exist
integers $x_1,x_2,x_3,x_4$ such that
\begin{equation}
x_1^3 + x_2^3 + x_3^3 + x_4^3 = n.
\label{eq:main_problem}
\end{equation}
Unlike Waring's problem for cubes, the variables in
\eqref{eq:main_problem} are allowed to be negative.  This changes the
character of the problem: congruence obstructions are weaker, but finding
explicit representations in the difficult residue classes remains a
computational and arithmetic challenge.

\subsection{Related work}
\label{subsec:related_work}

\subsubsection{Waring's problem}
Waring's problem concerns representing positive integers as sums of
nonnegative $k$-th powers.  For cubes, it is known that
$g(3)=9$~\cite{HardyWright2008}, meaning every positive integer can be
represented as the sum of nine nonnegative cubes.  This differs
substantially from \eqref{eq:main_problem}, where negative cubes are
permitted.

\subsubsection{Sum of three cubes}
The sum of three cubes problem asks for integers $x,y,z$ such that
$x^3+y^3+z^3=n$.  Its congruence obstruction modulo $9$ leaves only the
classes $n\not\equiv \pm 4\pmod 9$ as possible, and even within those
classes the search can be very hard.  Recent large-scale computations
resolved notable cases such as $33$ and $42$~\cite{Booker2019,
BookerSutherland2021}.

\subsubsection{Sum of five cubes}
The sum of five cubes problem is much easier.  The identity
\cite{HardyWright2008}
\begin{equation}
6x = (x+1)^3 + (x-1)^3 - x^3 - x^3,
\end{equation}
shows that any multiple of $6$ is a sum of four integer cubes.  Since
$r^3\bmod 6$ covers all residue classes, one additional cube is enough to
represent any integer.

\subsection{Known results and contribution}
\label{subsec:known_results}

In 1966, Demjanenko~\cite{Demjanenko1966} proved that every integer
$n\not\equiv 4,5\pmod 9$ is a sum of four integer cubes.  The remaining
classes $n\equiv \pm 4\pmod 9$ are the natural target for computation.
Demjanenko's identities are constructive, but in some cases they produce
large representations.

The present paper gives a direct computational verification for the full
range $1\leq n\leq 10^9$, including the classes $n\equiv\pm4\pmod 9$.
We do not claim that transforming four-cube identities into Pell-type
equations is new in itself; related substitutions are part of the
computational folklore around sums of cubes.  The contribution here is
instead threefold.  First, we isolate a simple one-parameter family with
explicit integrality, parity, and local congruence constraints.  Second, we
use this family to produce certificates for every $1\leq n\leq 10^9$.
Third, we provide an independently checkable verification protocol and an
empirical local-admissibility model for the observed search behavior.

\section{Preliminaries}
\label{sec:preliminaries}

\subsection{Generalized Pell equations}
\label{subsec:pell_equation}

A generalized Pell equation is a Diophantine equation of the form
$x^2 - dy^2 = m$, where $d$ is a non-square positive integer and $m$ is an
integer. When $m=1$, it is called the standard Pell equation.

There are mature algorithms for solving generalized Pell equations (e.g.,
PQa algorithm~\cite{Robertson2004}, LMM algorithm~\cite{Matthews2000}).
These algorithms have wide applications in computational number theory and
have efficient implementations. For example, the \texttt{solve\_integer}
method in SageMath~\cite{SageMath2023} is based on the PARI/GP
library~\cite{PARI2024} and Simon's reduction
algorithm~\cite{Simon2005}.

\subsection{Demjanenko's Theorem}
\label{subsec:demjanenko_theorem}

\begin{theorem}[Demjanenko, 1966~\cite{Demjanenko1966}]
\label{thm:demjanenko}
For any integer $n$ satisfying $n \not\equiv 4,5 \pmod{9}$, there exist integers $x_1,x_2,x_3,x_4$ such that $x_1^3 + x_2^3 + x_3^3 + x_4^3 = n$.
\end{theorem}

\section{Main Results}
\label{sec:main_results}

\subsection{Verified range}
\label{subsec:verified_range}

\begin{theorem}[Computational verification]
\label{thm:verification}
For every integer $n$ with $1\leq n\leq 10^9$, the computation described
in Section~\ref{sec:experimental} produced integers
$x_1,x_2,x_3,x_4$ satisfying
\[
        x_1^3+x_2^3+x_3^3+x_4^3=n.
\]
\end{theorem}

\begin{proof}
This is a computational proof.  Each recorded quadruple was verified by
exact integer arithmetic; the public data set records the representations,
and the verification protocol is described in
Section~\ref{subsec:verification_protocol}.
\end{proof}

\subsection{Transformation to Pell equation}
\label{subsec:transformation}

\begin{theorem}[Pell transformation with congruence constraints]
\label{thm:transformation}
Let $n\in\mathbb Z$ and let $d$ be a positive integer.  Put
\[
        M_{n,d}=\frac{4n-1+d^3}{3}.
\]
If $M_{n,d}\in\mathbb Z$, then the map
\[
        (x,y)\longmapsto (u,v)=(2x+1,2y+d)
\]
is a bijection between integer solutions of
\begin{equation}
(x+1)^3 - x^3 - (y+d)^3 + y^3 = n
\label{eq:original_form}
\end{equation}
and integer solutions of
\begin{equation}
u^2 - d v^2 = M_{n,d}
\label{eq:pell_form}
\end{equation}
satisfying
\[
        u\equiv 1\pmod 2,\qquad v\equiv d\pmod 2.
\]
Moreover, $M_{n,d}\in\mathbb Z$ if and only if
$d\equiv 1-n\pmod 3$.
\end{theorem}

\begin{proof}
Expanding the left-hand side of equation~\eqref{eq:original_form}:
\begin{align*}
&(x+1)^3 - x^3 - (y+d)^3 + y^3 \\
&= (x^3 + 3x^2 + 3x + 1) - x^3
     - (y^3 + 3d y^2 + 3d^2y + d^3) + y^3 \\
&= 3x^2 + 3x + 1 - 3d y^2 - 3d^2y - d^3.
\end{align*}

Rearranging terms:
\begin{equation}
3x^2 + 3x - 3d y^2 - 3d^2y = n - 1 + d^3.
\end{equation}

Multiplying both sides by 4 and completing the squares:
\begin{align*}
&12x^2 + 12x - 12d y^2 - 12d^2y
       = 4n - 4 + 4d^3 \\
&3(4x^2 + 4x) - 3d(4y^2 + 4d y)
       = 4n - 4 + 4d^3 \\
&3[(2x+1)^2 - 1] - 3d[(2y+d)^2 - d^2]
       = 4n - 4 + 4d^3 \\
&(2x+1)^2 - 1 - d(2y+d)^2 + d^3
       = \frac{4n - 4 + 4d^3}{3} \\
&u^2 - d v^2 = \frac{4n - 1 + d^3}{3}.
\end{align*}
Thus any integer pair $(x,y)$ gives a solution of
\eqref{eq:pell_form} with $u$ odd and $v\equiv d\pmod2$.
Conversely, if $(u,v)$ satisfies \eqref{eq:pell_form} and these two
congruences, then
\[
        x=\frac{u-1}{2},\qquad y=\frac{v-d}{2}
\]
are integers, and substituting them in the displayed calculation gives
\eqref{eq:original_form}.  Finally,
$4n-1+d^3\equiv n-1+d\pmod3$, so
$M_{n,d}$ is integral exactly when
$d\equiv 1-n\pmod3$.
\end{proof}

The original identity also gives the parity obstruction
$n\equiv 1+d\pmod2$.  Combining this with the integrality condition
shows that it is enough to search
\[
        d\equiv 1-n\pmod6.
\]

\begin{lemma}[Local obstructions]
\label{lem:local_obstruction}
Let $p^\alpha\Vert d$.  If
\[
        u^2-d v^2=M_{n,d}
\]
has an integer solution, then $M_{n,d}$ is a square modulo
$p^\alpha$.  In particular, for odd primes $p\neq 3$,
\[
        M_{n,d}\equiv (4n-1)\cdot 3^{-1}\pmod {p^\alpha},
\]
so $(4n-1)\cdot 3^{-1}$ must be a square modulo $p^\alpha$.
\end{lemma}

\begin{proof}
Reducing the equation modulo $p^\alpha$ gives
$u^2\equiv M_{n,d}\pmod {p^\alpha}$, since
$p^\alpha$ divides $d$.  If $p\neq3$, then $3$ is invertible modulo
$p^\alpha$, and $d^3\equiv0\pmod {p^\alpha}$, giving the stated
congruence.
\end{proof}

\subsection{Algorithm Description}
\label{subsec:algorithm}

\subsubsection{Algorithm Pseudocode}
\label{subsubsec:pseudocode}

\begin{algorithm}
\caption{SolveFourCubes}
\label{alg:main}
\begin{algorithmic}[1]
\Require Integer $n$
\Ensure Integers $X_1,X_2,X_3,X_4$ such that
$X_1^3+X_2^3+X_3^3+X_4^3=n$, or $\emptyset$ if no solution found
\If{$n \bmod 9 = 4$}
    \State $(X_1,X_2,X_3,X_4) \gets \textsc{SolveFourCubes}(-n)$
    \State \Return $(-X_1,-X_2,-X_3,-X_4)$
\EndIf
\For{positive integers $d\equiv 1-n\pmod6$}
    \If{$d$ is a perfect square}
        \State \textbf{continue}
    \EndIf
    \State $M \gets (4n - 1 + d^3)/3$
    \If{$M$ fails the local tests of Lemma~\ref{lem:local_obstruction}}
        \State \textbf{continue}
    \EndIf
    \State Let $\mathcal C$ be the quadruples obtained from solutions of
           $u^2-d v^2=M$ subject to
           $u\equiv1\pmod2$ and $v\equiv d\pmod2$
    \If{$\mathcal C=\emptyset$}
        \State \textbf{continue}
    \EndIf
    \State \Return the element of $\mathcal C$ minimizing
           $H=\max_i |X_i|$, with lexicographic tie-breaking
\EndFor
\State \Return $\emptyset$
\end{algorithmic}
\end{algorithm}

The congruence-constrained Pell step in Algorithm~\ref{alg:main} means that
we do not rely on a single arbitrary solution returned by a generalized Pell
solver.  The reference implementation enumerates all representatives returned
by the underlying quadratic-form solver, keeps those satisfying
$u\equiv1\pmod2$ and $v\equiv d\pmod2$, converts them to four-cube
quadruples, and chooses the one with minimal height
$H=\max_i|X_i|$.  The independent verification pass described in
Section~\ref{subsec:verification_protocol} does not depend on how this
solution was found.

\subsubsection{Handling Obstructions}
\label{subsubsec:obstructions}

The expression
$(x+1)^3-x^3-(y+d)^3+y^3$ cannot be congruent to $4$ modulo $9$,
but it can be congruent to $5$.  Therefore, when $n\equiv4\pmod9$ we
apply the same search to $-n$ and then negate the four resulting cubes.

\subsection{Why we use an indefinite Pell family}
\label{subsec:necessity}

The identity above is one member of a more general two-difference family
\begin{equation}
(x+a)^3 - x^3 + (y+b)^3 - y^3 = n,
\end{equation}
where $a,b$ are integers.  Expanding and completing squares gives an
equation of the form
\begin{equation}
a u^2 + b v^2 = m,
\end{equation}
where $u=2x+a$, $v=2y+b$, and
$m=(4n-a^3-b^3)/3$.  Positive-definite choices of $a$ and $b$ can produce
small representations in special cases, but they do not lead to Pell-type
unit orbits.  We use the mixed-sign family $a=1$, $b=-d$ because it
produces the indefinite norm equation $u^2-d v^2=M_{n,d}$, for
which generalized Pell solvers and congruence filtering are available.

\section{Local-admissibility heuristic}
\label{sec:theoretical_analysis}

\subsection{Local admissibility}
\label{subsec:local_admissibility}

The values $M_{n,d}=(4n-1+d^3)/3$ are not random integers
independent of $d$.  In particular, Lemma~\ref{lem:local_obstruction}
shows that a prime dividing $d$ can force the success probability for
that $d$ to be zero.  We therefore separate the local restrictions from
the global representation question.

For fixed $n$, define $A_n(d)=1$ if the following conditions hold:
\[
        d>0,\qquad d\ \text{is nonsquare},\qquad
        d\equiv 1-n\pmod6,
\]
and, for every prime power $p^\alpha\Vert d$, the integer
$M_{n,d}$ is a square modulo $p^\alpha$.  Otherwise set
$A_n(d)=0$.  Only the values with $A_n(d)=1$ are locally
admissible candidates for the Pell search.

For odd primes $p\neq3$ with $p\nmid 4n-1$, the local condition at a prime
factor $p$ of $d$ is controlled by the quadratic character
\[
        \chi_n(p)=
        \left(\frac{(4n-1)\cdot 3^{-1}}{p}\right).
\]
Thus a prime $p$ with $\chi_n(p)=-1$ cannot divide a locally admissible
$d$.  For typical $n$, and away from the finitely many primes dividing
$6(4n-1)$, these characters are modeled as taking the values $\pm1$ with
roughly equal frequency.

\subsection{Counting locally admissible parameters}
\label{subsec:counting_admissible}

The preceding character model suggests that the prime factors of a locally
admissible $d$ must lie in a set of primes of natural density about
$1/2$.  Standard sieve heuristics for integers supported on such a set then
predict
\[
        \#\{d\leq t:A_n(d)=1\}
        \approx C_n \frac{t}{\sqrt{\log t}},
\]
where $C_n$ depends on the exceptional local behavior at small primes and at
primes dividing $4n-1$.  This is the point at which the dependence of
$M_{n,d}$ on $d$ enters the model: locally impossible values of
$d$ are removed before any Bernays-type density is invoked.

\subsection{Conditional Bernays-type density}
\label{subsec:conditional_bernays}

Bernays' theorem for indefinite binary quadratic forms~\cite{Bernays1912}
motivates a global density of order $1/\sqrt{\log X}$ for integers up to
$X$ represented by a fixed indefinite form, after local restrictions are
taken into account.  We use this only as a conditional heuristic.  Given
$A_n(d)=1$, model the probability that
\[
        u^2-d v^2=M_{n,d}
\]
has a solution in the required parity class by
\[
        \Pr(\text{success at }d\mid A_n(d)=1)
        \approx \frac{K_n(d)}{\sqrt{\log M_{n,d}}},
\]
where $K_n(d)$ is treated as a slowly varying factor.  During the
early part of the search, when $d^3$ is small compared with $n$, we
have $\log M_{n,d}\asymp\log n$.

The expected number of successful candidates up to $t$ is then modeled by
\[
        \Lambda_n(t)
        \approx
        \sum_{\substack{d\leq t\\ A_n(d)=1}}
        \frac{K_n(d)}{\sqrt{\log M_{n,d}}}
        \approx
        \frac{C_nK_n\,t}{\sqrt{\log t}\sqrt{\log n}},
\]
again up to $n$-dependent and slowly varying constants.  Solving
$\Lambda_n(t)\approx1$ suggests that the first successful parameter should
typically have size
\[
        t \approx \sqrt{\log n\,\log\log n},
\]
within these constants.  We do not state this as a theorem; it is a
local-admissibility stopping-time model, used to interpret the empirical
distribution in Section~\ref{subsec:distribution}.

\section{Experimental Verification}
\label{sec:experimental}

\subsection{Verification protocol}
\label{subsec:verification_protocol}

For each integer $n$ in the range $1\leq n\leq 10^9$, the program searches
over positive nonsquare integers $d$ as in Algorithm~\ref{alg:main}.
If $n\equiv 4\pmod 9$, the program applies the same search to $-n$ and
then negates the resulting four cubes.  Once a candidate quadruple
$(x_1,x_2,x_3,x_4)$ is found, it is not accepted until the equality
\[
        x_1^3+x_2^3+x_3^3+x_4^3=n
\]
has been checked using exact integer arithmetic.

At the level needed for independent verification, a public certificate row
records $n$ and the four integers $x_1,x_2,x_3,x_4$.  The search parameter
$d$ used in the tables is recoverable from the generated family as
$d=|x_3+x_4|$, but the final checker does not rely on this value.  The
checker verifies that the records cover the interval without gaps or
repetitions and that each recorded quadruple satisfies the defining equation.
No floating-point arithmetic and no Pell solver are used in this final check.

\subsection{Certificate and independent checking}
\label{subsec:certificate}

The computational proof is organized so that the expensive search and the
final verification are logically separated.  The generator may use
generalized Pell solvers, congruence filters, and implementation-specific
optimizations.  The checker treats the output as a list of certificates and
performs only elementary integer tests:
\[
        1\leq n\leq 10^9,\qquad
        x_1^3+x_2^3+x_3^3+x_4^3=n.
\]
It also checks that the values of $n$ occur exactly once.  Thus a failure
or incompleteness in the Pell search would be detected by the independent
checking pass, and the correctness of the final certificate does not depend
on the heuristic model of Section~\ref{sec:theoretical_analysis}.

\subsection{Results}
\label{subsec:results}

Using this protocol, we generated a public certificate data set for every
integer $0\leq n\leq 10^9$, stored in $100$ shards.  The first $99$ shards
have length $10^7$, and the final shard contains
$990000000\leq n\leq 10^9$.  Restricting to $1\leq n\leq 10^9$ gives
Theorem~\ref{thm:verification}.  Solutions for negative integers $-n$ are
obtained by taking the negatives of the four entries.  Some small examples
are listed in Table~\ref{tab:solutions}.

\subsection{Data availability}
The associated code and data are publicly available on Hugging Face at
\url{https://huggingface.co/datasets/}\texttt{kzoacn/sum4cubes}.  The files
\texttt{cubes-00.txt},\ldots,\texttt{cubes-99.txt} cover
$0\leq n\leq 10^9$.

\subsection{Distribution analysis}
\label{subsec:distribution}

Due to computational time constraints, we performed detailed statistical
analysis for the reference implementation only on the range
$1 \leq n \leq 10^7$.  In the plots below, $d(n)$ denotes the first search
parameter $d$ for which the reference implementation finds at least one
parity-compatible Pell representative.  For that $d$, the displayed
quadruple is the one of minimal height $H(n)=\max_i |x_i|$, with
lexicographic tie-breaking.  The figures were regenerated from this rule
using bins of length $10^4$ and a deterministic sample of every $5000$-th
integer.

Figure~\ref{fig:points} shows sampled values of $d(n)$ on a logarithmic
scale, with the largest outliers in this range highlighted.  Most values of
$d(n)$ are small, while the largest values are separated by several orders
of magnitude.

Figure~\ref{fig:d_quantiles} gives binned quantile curves for
$d(n)$.  The median and high quantiles remain stable as $n$ grows, while
the bin maxima record rare large deviations.

Figure~\ref{fig:average} shows the cumulative average of $d(n)$ together
with a least-squares fit to the heuristic shape
$a\sqrt{\log n\log\log n}+b$.

For the full range $1\leq n\leq 10^9$, Appendix
Tables~\ref{tab:largest_d_1e9} and~\ref{tab:largest_height_1e9} summarize the
regenerated certificate data using the same rule for $d(n)$ and $H(n)$.
The leading examples in these two senses are written out explicitly in
Tables~\ref{tab:rep_936384844} and~\ref{tab:rep_340984507}.

\begin{figure}[H]
\centering
\begin{tikzpicture}
\begin{axis}[
    width=0.98\linewidth,
    height=0.58\linewidth,
    ymode=log,
    xmin=0,
    xmax=10000000,
    ymin=1,
    ymax=1000000,
    xlabel={$n$},
    ylabel={$d(n)$},
    scaled x ticks=false,
    xtick={0,2000000,4000000,6000000,8000000,10000000},
    xticklabels={$0$,$2\cdot10^6$,$4\cdot10^6$,$6\cdot10^6$,$8\cdot10^6$,$10^7$},
    grid=both,
    major grid style={black!12},
    minor grid style={black!6},
    tick label style={font=\small},
    label style={font=\small},
    legend columns=2,
    legend style={draw=none, fill=none, font=\small, at={(0.5,1.02)}, anchor=south},
    legend cell align=left,
]
\addplot[
    only marks,
    mark=*,
    mark size=0.35pt,
    draw=none,
    fill=black,
    fill opacity=0.16,
] table[x=n,y=d] {figures/data/d_sample.dat};
\addlegendentry{sample}

\addplot[
    only marks,
    mark=*,
    mark size=1.7pt,
    draw=red!70!black,
    fill=red!70!black,
] table[x=n,y=d] {figures/data/d_outliers.dat};
\addlegendentry{top ten}
\end{axis}
\end{tikzpicture}
\caption{Reference search parameters $d(n)$ for $1\leq n\leq 10^7$ on a logarithmic scale.  The gray points are the deterministic sample $n\equiv 0\pmod {5000}$, and the red points mark the ten largest observed values of $d(n)$ in this range.}
\label{fig:points}
\end{figure}


\begin{figure}[H]
\centering
\begin{tikzpicture}
\begin{axis}[
    width=0.98\linewidth,
    height=0.58\linewidth,
    ymode=log,
    xmin=0,
    xmax=10000000,
    ymin=1,
    ymax=1000000,
    xlabel={$n$},
    ylabel={$d(n)$},
    scaled x ticks=false,
    xtick={0,2000000,4000000,6000000,8000000,10000000},
    xticklabels={$0$,$2\cdot10^6$,$4\cdot10^6$,$6\cdot10^6$,$8\cdot10^6$,$10^7$},
    grid=both,
    major grid style={black!12},
    minor grid style={black!6},
    tick label style={font=\small},
    label style={font=\small},
    legend columns=3,
    legend style={draw=none, fill=white, font=\small, at={(0.5,1.02)}, anchor=south},
    legend cell align=left,
]
\addplot[very thick, blue!70!black] table[x=n,y=median] {figures/data/d_quantiles.dat};
\addlegendentry{median}

\addplot[thick, teal!70!black] table[x=n,y=p90] {figures/data/d_quantiles.dat};
\addlegendentry{$90\%$}

\addplot[thick, orange!85!black] table[x=n,y=p99] {figures/data/d_quantiles.dat};
\addlegendentry{$99\%$}

\addplot[thick, red!70!black] table[x=n,y=p999] {figures/data/d_quantiles.dat};
\addlegendentry{$99.9\%$}

\addplot[thin, black!60] table[x=n,y=max] {figures/data/d_quantiles.dat};
\addlegendentry{bin max}
\end{axis}
\end{tikzpicture}
\caption{Binned quantiles of the reference search parameter $d(n)$ for $1\leq n\leq 10^7$.  Each bin contains $10^4$ consecutive values of $n$.}
\label{fig:d_quantiles}
\end{figure}


\begin{figure}[H]
\centering
\begin{tikzpicture}
\begin{axis}[
    width=0.98\linewidth,
    height=0.58\linewidth,
    xmin=0,
    xmax=10000000,
    ymin=25,
    ymax=46,
    xlabel={$N$},
    ylabel={$\frac{1}{N}\sum_{n\leq N}d(n)$},
    scaled x ticks=false,
    xtick={0,2000000,4000000,6000000,8000000,10000000},
    xticklabels={$0$,$2\cdot10^6$,$4\cdot10^6$,$6\cdot10^6$,$8\cdot10^6$,$10^7$},
    grid=both,
    major grid style={black!12},
    minor grid style={black!6},
    tick label style={font=\small},
    label style={font=\small},
    legend columns=2,
    legend style={draw=none, fill=none, font=\small, at={(0.5,1.02)}, anchor=south},
    legend cell align=left,
]
\addplot[very thick, blue!70!black] table[x=n,y=avg] {figures/data/d_average.dat};
\addlegendentry{observed average}

\addplot[thick, dashed, red!70!black] table[x=n,y=model] {figures/data/d_average.dat};
\addlegendentry{least-squares fit}
\end{axis}
\end{tikzpicture}
\caption{Cumulative average of the reference search parameter $d(n)$ for $1\leq n\leq 10^7$, compared with a least-squares fit of the form $a\sqrt{\log N\log\log N}+b$.}
\label{fig:average}
\end{figure}


\subsection{Running Time}
\label{subsec:running_time}

The algorithm also runs quickly for larger integers, e.g., $n = 9 \cdot 2^{128} + 5$:
\begin{align*}
n &= 3062541302288446171170371466885913903109 \\
&= 39512676350251656797^3 - 39512676350251656796^3 \\
&\quad - 1275823017616426120^3 + 1275823017616425788^3.
\end{align*}
This solution was found within seconds.
Appendix Tables~\ref{tab:rep_pi_prefix} and~\ref{tab:rep_e_prefix} give
two further examples obtained from decimal prefixes of \(\pi\) and \(e\),
with the decimal point omitted and a prefix length chosen so that
\(n\equiv 4\pmod 9\).

\subsection{Comparison with Demjanenko's Method}
\label{subsec:comparison}

For $n = 254$, our algorithm yields:
\begin{equation}
(-12)^3 + 13^3 + 1^3 + (-6)^3 = 254,
\end{equation}
while Demjanenko's method~\cite{Demjanenko1966,Alpertron2020} would classify it as a special case modulo 108, outputting four numbers with over 600 digits each (see \url{https://www.alpertron.com.ar/FCUBES.HTM}).

\section*{Acknowledgments}

We thank the anonymous reviewers for their careful reading and constructive
suggestions.  In particular, we are grateful for the observation that a single
generalized Pell equation may have both parity-compatible and
parity-incompatible solutions, which led us to state the congruence constraints
explicitly and to use a congruence-aware Pell search in the final
implementation.

\bibliographystyle{amsplain}
\bibliography{ref}

\appendix
\section{SageMath Code Implementation}
\label{app:code}

The following executable SageMath code records a reference implementation of
the search logic.  It was tested with SageMath 10.7.  The production
implementation uses the same certificate format, while the independent checker
described in Section~\ref{subsec:certificate} verifies only the resulting
four-cube identities.  The same code is included in the source tree as
\texttt{scripts/four\_cubes\_reference.py}.

\begin{verbatim}
from sage.all import ZZ, BinaryQF, Mod, factor, is_square


def local_admissible(n, d, M):
    """
    Necessary local tests from Lemma 3.3.
    For each p^a || d, M must be a square modulo p^a.
    """
    for p, a in factor(d):
        q = p**a
        if not Mod(M, q).is_square():
            return False
    return True


def solve_pell_with_congruences(d, M):
    """
    Yield solutions of
        u^2 - d*v^2 = M,
        u = 1 (mod 2), v = d (mod 2),
    returned by Sage in this congruence class.
    """
    Q = BinaryQF([1, 0, -d])
    for u, v in Q.solve_integer(M, _flag=3):
        u, v = ZZ(u), ZZ(v)
        if (u - 1) % 2 == 0 and (v - d) % 2 == 0:
            yield u, v


def quadruple_from_pell(d, u, v):
    x = (u - 1) // 2
    y = (v - d) // 2
    return (x + 1, -x, -(y + d), y)


def height(xs):
    return max(abs(x) for x in xs)


def solve(n, LIMIT=10**8, return_d=False):
    """
    Solve x1^3 + x2^3 + x3^3 + x4^3 = n by the
    congruence-filtered Pell family.
    """
    n = ZZ(n)
    if n % 9 == 4:
        sol = solve(-n, LIMIT, return_d=True)
        if sol is None:
            return None
        x1, x2, x3, x4, d = sol
        out = (-x1, -x2, -x3, -x4)
        return out + (d,) if return_d else out

    residue = ZZ((1 - n) % 6)
    if residue == 0:
        residue = ZZ(6)

    for d in range(residue, LIMIT + 1, 6):
        d = ZZ(d)
        if is_square(d):
            continue

        numerator = 4*n - 1 + d**3
        if numerator % 3 != 0:
            continue
        M = numerator // 3

        if not local_admissible(n, d, M):
            continue

        candidates = [
            quadruple_from_pell(d, u, v)
            for u, v in solve_pell_with_congruences(d, M)
        ]
        if not candidates:
            continue

        out = min(
            candidates,
            key=lambda xs: (height(xs), tuple(xs)),
        )
        return out + (d,) if return_d else out

    return None


def verify(n, xs):
    return sum(ZZ(x)**3 for x in xs) == ZZ(n)


for n in [254, 314159265358979323, 27182818284590452]:
    sol = solve(n, LIMIT=2000, return_d=True)
    assert sol is not None
    assert verify(n, sol[:4])
    print(n, sol)
\end{verbatim}

\section{Tables}
\label{app:tables}

\newcolumntype{N}{>{\raggedright\arraybackslash\ttfamily\footnotesize}p{0.78\textwidth}}

For a representation produced by the reference implementation in
Algorithm~\ref{alg:main}, we write $d(n)$ for the first search parameter
found by the algorithm and
\[
        H(n)=\max_{1\leq i\leq 4}|x_i|
\]
for the height of the corresponding minimal-height representative at that
search parameter.  The first table below lists small examples explicitly.
The following two tables select outlying instances from the regenerated
certificate data for the full verification range $1\leq n\leq 10^9$.
The column $\ell(H(n))$ denotes the number of decimal digits of $H(n)$.

\begin{longtable}{c c c}
\caption{Sum of four cubes representations for $1\leq n \leq 100$} \label{tab:solutions}\\
\toprule
$n$ & $x_1^3 + x_2^3 + x_3^3 + x_4^3$ & $d$ \\
\midrule
\endfirsthead
\toprule
$n$ & $x_1^3 + x_2^3 + x_3^3 + x_4^3$ & $d$ \\
\midrule
\endhead
\midrule
\multicolumn{3}{r}{Continued on next page} \\
\endfoot
\bottomrule
\endlastfoot
1 & $(-6)^3 + 7^3 + (-1)^3 + (-5)^3$ & 6 \\
2 & $(-3)^3 + 4^3 + (-3)^3 + (-2)^3$ & 5 \\
3 & $235^3 + (-234)^3 + (-56)^3 + 22^3$ & 34 \\
4 & $4^3 + (-5)^3 + 4^3 + 1^3$ & 5 \\
5 & $2^3 + (-1)^3 + (-1)^3 + (-1)^3$ & 2 \\
6 & $(-58)^3 + 59^3 + (-21)^3 + (-10)^3$ & 31 \\
7 & $5^3 + (-4)^3 + (-3)^3 + (-3)^3$ & 6 \\
8 & $14^3 + (-13)^3 + (-8)^3 + (-3)^3$ & 11 \\
9 & $(-108)^3 + 109^3 + (-27)^3 + (-25)^3$ & 52 \\
10 & $(-2)^3 + 3^3 + (-2)^3 + (-1)^3$ & 3 \\
11 & $(-2)^3 + 3^3 + (-2)^3 + 0^3$ & 2 \\
12 & $(-186)^3 + 187^3 + (-47)^3 + (-8)^3$ & 55 \\
13 & $10^3 + (-11)^3 + 7^3 + 1^3$ & 8 \\
14 & $(-235)^3 + 236^3 + 2^3 + (-55)^3$ & 53 \\
15 & $842^3 + (-841)^3 + (-132)^3 + 56^3$ & 76 \\
16 & $(-19)^3 + 20^3 + (-10)^3 + (-5)^3$ & 15 \\
17 & $3^3 + (-2)^3 + (-1)^3 + (-1)^3$ & 2 \\
18 & $(-95)^3 + 96^3 + (-7)^3 + (-30)^3$ & 37 \\
19 & $(-5)^3 + 6^3 + (-4)^3 + (-2)^3$ & 6 \\
20 & $(-20)^3 + 21^3 + (-9)^3 + (-8)^3$ & 17 \\
21 & $10^3 + (-9)^3 + (-5)^3 + (-5)^3$ & 10 \\
22 & $85^3 + (-86)^3 + 28^3 + 1^3$ & 29 \\
23 & $146^3 + (-145)^3 + (-40)^3 + 8^3$ & 32 \\
24 & $(-16)^3 + 17^3 + (-9)^3 + (-4)^3$ & 13 \\
25 & $(-136)^3 + 137^3 + (-38)^3 + (-10)^3$ & 48 \\
26 & $(-4)^3 + 5^3 + (-3)^3 + (-2)^3$ & 5 \\
27 & $(-11)^3 + 12^3 + (-7)^3 + (-3)^3$ & 10 \\
28 & $(-3)^3 + 4^3 + (-2)^3 + (-1)^3$ & 3 \\
29 & $(-3)^3 + 4^3 + (-2)^3 + 0^3$ & 2 \\
30 & $120^3 + (-119)^3 + (-35)^3 + 4^3$ & 31 \\
31 & $229687^3 + (-229688)^3 + 7390^3 + (-6260)^3$ & 1130 \\
32 & $431^3 + (-430)^3 + (-94)^3 + 65^3$ & 29 \\
33 & $59^3 + (-58)^3 + (-19)^3 + (-15)^3$ & 34 \\
34 & $(-4)^3 + 5^3 + (-3)^3 + 0^3$ & 3 \\
35 & $4^3 + (-3)^3 + (-1)^3 + (-1)^3$ & 2 \\
36 & $(-6)^3 + 7^3 + (-4)^3 + (-3)^3$ & 7 \\
37 & $6^3 + (-5)^3 + (-3)^3 + (-3)^3$ & 6 \\
38 & $(-41)^3 + 42^3 + (-17)^3 + (-6)^3$ & 23 \\
39 & $117^3 + (-116)^3 + (-35)^3 + 13^3$ & 22 \\
40 & $148^3 + (-149)^3 + 40^3 + 13^3$ & 53 \\
41 & $8^3 + (-7)^3 + (-4)^3 + (-4)^3$ & 8 \\
42 & $(-109)^3 + 110^3 + 15^3 + (-34)^3$ & 19 \\
43 & $(-7)^3 + 8^3 + (-5)^3 + (-1)^3$ & 6 \\
44 & $(-7)^3 + 8^3 + (-5)^3 + 0^3$ & 5 \\
45 & $(-44)^3 + 45^3 + (-18)^3 + (-4)^3$ & 22 \\
46 & $25^3 + (-24)^3 + (-3)^3 + (-12)^3$ & 15 \\
47 & $(-9)^3 + 10^3 + (-6)^3 + (-2)^3$ & 8 \\
48 & $(-21)^3 + 22^3 + (-11)^3 + (-2)^3$ & 13 \\
49 & $(-2)^3 + 1^3 + 4^3 + (-2)^3$ & 2 \\
50 & $(-76)^3 + 77^3 + (-22)^3 + (-19)^3$ & 41 \\
51 & $(-10)^3 + 11^3 + (-6)^3 + (-4)^3$ & 10 \\
52 & $(-4)^3 + 5^3 + (-2)^3 + (-1)^3$ & 3 \\
53 & $(-4)^3 + 5^3 + (-2)^3 + 0^3$ & 2 \\
54 & $(-9)^3 + 10^3 + (-6)^3 + (-1)^3$ & 7 \\
55 & $(-6)^3 + 7^3 + (-4)^3 + (-2)^3$ & 6 \\
56 & $(-5)^3 + 6^3 + (-3)^3 + (-2)^3$ & 5 \\
57 & $(-21)^3 + 22^3 + 1^3 + (-11)^3$ & 10 \\
58 & $1^3 + (-2)^3 + 4^3 + 1^3$ & 5 \\
59 & $5^3 + (-4)^3 + (-1)^3 + (-1)^3$ & 2 \\
60 & $1202^3 + (-1201)^3 + (-163)^3 + 0^3$ & 163 \\
61 & $(-16)^3 + 17^3 + (-9)^3 + (-3)^3$ & 12 \\
62 & $7^3 + (-6)^3 + (-4)^3 + (-1)^3$ & 5 \\
63 & $(-30)^3 + 31^3 + (-12)^3 + (-10)^3$ & 22 \\
64 & $(-5)^3 + 6^3 + (-3)^3 + 0^3$ & 3 \\
65 & $(-5)^3 + 6^3 + (-3)^3 + 1^3$ & 2 \\
66 & $(-198)^3 + 199^3 + (-50)^3 + 19^3$ & 31 \\
67 & $(-5)^3 + 4^3 + 4^3 + 4^3$ & 8 \\
68 & $257^3 + (-256)^3 + (-58)^3 + (-13)^3$ & 71 \\
69 & $(-148)^3 + 149^3 + (-42)^3 + 20^3$ & 22 \\
70 & $(-58)^3 + 59^3 + (-20)^3 + (-13)^3$ & 33 \\
71 & $(-6)^3 + 7^3 + 2^3 + (-4)^3$ & 2 \\
72 & $(-14)^3 + 15^3 + (-7)^3 + (-6)^3$ & 13 \\
73 & $7^3 + (-6)^3 + (-3)^3 + (-3)^3$ & 6 \\
74 & $25^3 + (-24)^3 + (-12)^3 + 1^3$ & 11 \\
75 & $(-53)^3 + 54^3 + (-20)^3 + (-8)^3$ & 28 \\
76 & $10^3 + (-11)^3 + 7^3 + 4^3$ & 11 \\
77 & $(-19)^3 + 20^3 + (-10)^3 + (-4)^3$ & 14 \\
78 & $(-7)^3 + 8^3 + (-4)^3 + (-3)^3$ & 7 \\
79 & $(-13)^3 + 14^3 + (-7)^3 + (-5)^3$ & 12 \\
80 & $(-16)^3 + 17^3 + (-9)^3 + (-2)^3$ & 11 \\
81 & $11^3 + (-10)^3 + (-5)^3 + (-5)^3$ & 10 \\
82 & $(-5)^3 + 6^3 + (-2)^3 + (-1)^3$ & 3 \\
83 & $(-5)^3 + 6^3 + (-2)^3 + 0^3$ & 2 \\
84 & $(-8)^3 + 9^3 + (-5)^3 + (-2)^3$ & 7 \\
85 & $(-86)^3 + 85^3 + 28^3 + 4^3$ & 32 \\
86 & $(-106795)^3 + 106796^3 + (-4039)^3 + 3164^3$ & 875 \\
87 & $(-16)^3 + 17^3 + (-1)^3 + (-9)^3$ & 10 \\
88 & $473^3 + (-472)^3 + (-90)^3 + 39^3$ & 51 \\
89 & $6^3 + (-5)^3 + (-1)^3 + (-1)^3$ & 2 \\
90 & $(-188)^3 + 189^3 + (-43)^3 + (-30)^3$ & 73 \\
91 & $9^3 + (-8)^3 + (-5)^3 + (-1)^3$ & 6 \\
92 & $(-6)^3 + 7^3 + (-3)^3 + (-2)^3$ & 5 \\
93 & $165^3 + (-164)^3 + (-44)^3 + 16^3$ & 28 \\
94 & $535^3 + (-536)^3 + 94^3 + 31^3$ & 125 \\
95 & $(-130)^3 + 131^3 + (-7)^3 + (-37)^3$ & 44 \\
96 & $17^3 + (-16)^3 + (-9)^3 + 2^3$ & 7 \\
97 & $(-7)^3 + 8^3 + (-4)^3 + (-2)^3$ & 6 \\
98 & $26^3 + (-25)^3 + (-12)^3 + (-5)^3$ & 17 \\
99 & $13^3 + (-12)^3 + (-7)^3 + (-3)^3$ & 10 \\
100 & $7^3 + (-6)^3 + (-3)^3 + 0^3$ & 3 \\

\end{longtable}

\begin{table}[H]
\centering
\caption{Largest recorded values of $d(n)$ for $1\leq n\leq 10^9$}
\label{tab:largest_d_1e9}
\begin{tabular}{c r r r r}
\toprule
Rank & $n$ & $d$ & $\ell(H)$ & $\log_{10}(H/n)$ \\
\midrule
1 & 936384844 & 5508491 & 66 & 56.516 \\
2 & 630252220 & 5170439 & 17 & 7.476 \\
3 & 172116347 & 3042818 & 26 & 17.569 \\
4 & 720368492 & 2483135 & 23 & 14.040 \\
5 & 62809898 & 2259611 & 18 & 9.591 \\
6 & 424484014 & 1890905 & 41 & 31.940 \\
7 & 984206885 & 1774796 & 27 & 17.274 \\
8 & 233143708 & 1759577 & 16 & 6.785 \\
9 & 200441272 & 1539353 & 33 & 24.389 \\
10 & 699315232 & 1421681 & 14 & 4.933 \\
\bottomrule
\end{tabular}
\end{table}

\begin{table}[H]
\centering
\caption{Largest recorded values of $H(n)$ for $1\leq n\leq 10^9$}
\label{tab:largest_height_1e9}
\begin{tabular}{c r r r r}
\toprule
Rank & $n$ & $d$ & $\ell(H)$ & $\log_{10}(H/n)$ \\
\midrule
1 & 340984507 & 1248734 & 131 & 122.036 \\
2 & 936384844 & 5508491 & 66 & 56.516 \\
3 & 139360442 & 473225 & 58 & 49.605 \\
4 & 812091910 & 1387697 & 57 & 48.084 \\
5 & 298597495 & 908216 & 51 & 41.866 \\
6 & 243310568 & 676853 & 47 & 38.513 \\
7 & 28875326 & 218663 & 44 & 36.337 \\
8 & 424484014 & 1890905 & 41 & 31.940 \\
9 & 584188196 & 94217 & 35 & 25.611 \\
10 & 200441272 & 1539353 & 33 & 24.389 \\
\bottomrule
\end{tabular}
\end{table}

\subsection{Explicit representations for two outliers}

The following two tables give the complete representations for the leading
outliers in Tables~\ref{tab:largest_d_1e9} and
\ref{tab:largest_height_1e9}.  In these tables, spaces in decimal
expansions are only digit separators, inserted in blocks of four for
readability.

\begin{table}[H]
\centering
\caption{Explicit certificate representation for $n=936384844$, the largest recorded value of $d(n)$}
\label{tab:rep_936384844}
\begin{tabular}{c N}
\toprule
Variable & Value \\
\midrule
$x_1$ & 30 7230 6200 3619 8887 6003 3192 9948 7199 7824 7989 0641 1712 7246 2307 8629 5941 \\
$x_2$ & -30 7230 6200 3619 8887 6003 3192 9948 7199 7824 7989 0641 1712 7246 2307 8629 5942 \\
$x_3$ & 130 9025 7064 0549 0465 2247 7828 1922 0428 4025 4310 9694 6153 1423 2834 8471 \\
$x_4$ & -130 9025 7064 0549 0465 2247 7828 1922 0428 4025 4310 9694 6153 1423 2283 9980 \\
\bottomrule
\end{tabular}
\[
\begin{aligned}
        936384844 &= x_1^3+x_2^3+x_3^3+x_4^3,\\
        d(936384844) &= 5508491,\qquad H(936384844)=|x_2|.
\end{aligned}
\]
\end{table}

\begin{table}[H]
\centering
\caption{Explicit certificate representation for $n=340984507$, the largest recorded value of $H(n)$}
\label{tab:rep_340984507}
\begin{tabular}{c N}
\toprule
Variable & Value \\
\midrule
$x_1$ & 370 5739 4133 4654 4780 1216 2360 0386 9785 7350 5408 0417 7220 5119 4147 2288 7378 2670 5430 4572 5176 0745 4543 2018 9726 4367 8218 0837 8821 0789 8314 1962 2346 \\
$x_2$ & -370 5739 4133 4654 4780 1216 2360 0386 9785 7350 5408 0417 7220 5119 4147 2288 7378 2670 5430 4572 5176 0745 4543 2018 9726 4367 8218 0837 8821 0789 8314 1962 2347 \\
$x_3$ & -3316 1938 4003 7246 7141 7215 0547 4292 0031 4035 9034 0369 9705 9982 1979 2976 7723 6077 8586 5629 2076 3931 3899 4614 3197 6791 5676 7171 1428 1184 4392 6711 \\
$x_4$ & 3316 1938 4003 7246 7141 7215 0547 4292 0031 4035 9034 0369 9705 9982 1979 2976 7723 6077 8586 5629 2076 3931 3899 4614 3197 6791 5676 7171 1428 1184 4517 5445 \\
\bottomrule
\end{tabular}
\[
\begin{aligned}
        340984507 &= x_1^3+x_2^3+x_3^3+x_4^3,\\
        d(340984507) &= 1248734,\qquad H(340984507)=|x_2|.
\end{aligned}
\]
\end{table}

\subsection{Decimal-prefix examples}

The following two tables give additional examples obtained from decimal
prefixes of \(\pi\) and \(e\).  The digit count includes the digit before the
decimal point; the decimal point is then omitted.  Both selected prefixes
satisfy \(n\equiv 4\pmod 9\).  The explicit representations are shown in
Tables~\ref{tab:rep_pi_prefix} and~\ref{tab:rep_e_prefix}.

\begin{table}[H]
\centering
\caption{Explicit representation for the 18-digit decimal prefix of \(\pi\)}
\label{tab:rep_pi_prefix}
\begin{tabular}{c@{\qquad}r}
\toprule
Variable & Value \\
\midrule
$x_1$ & 16411235899 \\
$x_2$ & -16411235900 \\
$x_3$ & 525952624 \\
$x_4$ & -525951650 \\
\bottomrule
\end{tabular}
\[
\begin{aligned}
        314159265358979323 &= x_1^3+x_2^3+x_3^3+x_4^3,\\
        d(314159265358979323) &= 974,\qquad H(314159265358979323)=|x_2|.
\end{aligned}
\]
\end{table}

\begin{table}[H]
\centering
\caption{Explicit representation for the 17-digit decimal prefix of \(e\)}
\label{tab:rep_e_prefix}
\begin{tabular}{c@{\qquad}r}
\toprule
Variable & Value \\
\midrule
$x_1$ & 6482848 \\
$x_2$ & -6482849 \\
$x_3$ & -19894238 \\
$x_4$ & 19894261 \\
\bottomrule
\end{tabular}
\[
\begin{aligned}
        27182818284590452 &= x_1^3+x_2^3+x_3^3+x_4^3,\\
        d(27182818284590452) &= 23,\qquad H(27182818284590452)=|x_4|.
\end{aligned}
\]
\end{table}

  

\end{document}
